\documentclass[11pt]{article}
\usepackage[a4paper,margin=1.1in]{geometry}
\usepackage[T1]{fontenc}
\usepackage{lmodern}
\usepackage{amsmath,amssymb,amsthm}
\usepackage{booktabs}
\usepackage{microtype}
\usepackage[hidelinks]{hyperref}
\usepackage{natbib}
\usepackage{setspace}
\usepackage{enumitem}

\newtheorem{theorem}{Theorem}
\newtheorem{proposition}[theorem]{Proposition}
\newtheorem{lemma}[theorem]{Lemma}
\newtheorem{corollary}[theorem]{Corollary}
\theoremstyle{definition}
\newtheorem{assumption}{Hypothesis}
\theoremstyle{remark}
\newtheorem{remark}[theorem]{Remark}

\newcommand{\leanunderscore}{\char`\_\hskip 0pt plus 0.4pt\relax}
\newcommand{\lean}[1]{\begingroup\ttfamily\small\let\_\leanunderscore #1\endgroup}
\newcommand{\formal}[1]{\par\noindent\begingroup\footnotesize\raggedright\emph{Formal statement:} \lean{#1}.\par\endgroup}
\newcommand{\R}{\mathbb{R}}
\newcommand{\dd}{\mathrm{d}}

\onehalfspacing
\tolerance=600 \emergencystretch=1.5em

\title{The Likelihood of a Computed Dynamic Model:\\ Proposition 2 of
Fern\'andez-Villaverde, Rubio-Ram\'irez and Santos, Corrected}
\author{LeanEconomics Org [mostly Claude]\thanks{The Lean development is
\texttt{LeanEconomics/Econometrics/ApproxLikelihood.lean} in
\url{https://github.com/LeanEconomics/LeanEconomics}. It was written mostly by Claude (Anthropic),
an AI coding model, working under the direction of Robert Kirkby (Victoria University of
Wellington); every result marked with a formal statement was checked by the Lean kernel and uses
only Lean's three standard axioms.}}
\date{October 2026\\[0.5em]\small Draft}

\begin{document}
\maketitle

\begin{abstract}
\noindent
Proposition 2 of \citet{FRS2006} bounds the gap between the exact and the computed likelihood of
a dynamic model by a term linear in the sample length, $T B\chi\delta$, with constants that do not
depend on $T$. \citet{AGH2009}, following \citet{Geweke2007}, show by a Gaussian counterexample
that this is false. We locate the error in a single step of the proof: a difference of two
$T$-fold products of densities is bounded as if the other $T-1$ factors were absent, so that the
``constant'' $B$ is really $M^{T-1}$, where $M$ bounds the density. We then show that the
proposition survives, with the same shape, once it is stated for the \emph{logarithm} of the
likelihood: under the paper's lower bound on densities, a pathwise version of its contraction
condition, and a Lipschitz condition on the log density in the initial state,
$|\log L_j(y^T) - \log L(y^T)| \le (T\chi/\xi + K/(1-\alpha))\,\delta$. The dependence on $T$ is
sharp. The average log-likelihood error is therefore $O(\delta)$ uniformly in $T$, which is the
quantity Ackerberg, Geweke and Hahn identify as relevant for estimation. The argument
is formalised in Lean~4 with Mathlib, including the existence of the coupling of invariant
distributions (by Krylov--Bogolyubov), the derivation of the hypotheses from primitives in the
case where the measurement shock alone is recovered from the data, and explicit bounds on the
resulting posteriors.
\end{abstract}

\section{Setting}

\citet{FRS2006} (FRS) study a state-space model
\[
  S_t = \varphi(S_{t-1}, W_t;\gamma), \qquad Y_t = g(S_t, V_t;\gamma),
\]
with compact supports, a unique invariant distribution $\mu^*$ of the state, and computed
approximations $\varphi_j, g_j$ satisfying $\|\varphi_j - \varphi\| \le \delta$ and
$\|g_j - g\| \le \delta$ in the sup norm. The shock $W_t$ is split as $(W_{1t}, W_{2t})$, where
$(W_{2t}, V_t)$ is recovered from the data by inverting the system. Write
$\omega = W_1^T$ for the unobserved shock path and
\[
  p_t(s,\omega) = p\bigl(y_t \mid W_1^t = \omega_{\le t},\, S_0 = s,\, y^{t-1};\gamma\bigr)
\]
for the conditional density of the $t$-th observation, with $q_t$ its counterpart under
$(\varphi_j, g_j)$. Then
\[
  L(y^T) = \int\!\!\int \prod_{t=1}^T p_t(s,\omega)\, Q(\dd\omega)\,\mu^*(\dd s), \qquad
  L_j(y^T) = \int\!\!\int \prod_{t=1}^T q_t(s,\omega)\, Q(\dd\omega)\,\mu_j^*(\dd s),
\]
where $Q$ is the law of $W_1^T$ and $\mu_j^*$ the invariant distribution of the computed chain.
FRS's Assumptions~4 and~8 give a lower bound $p_t, q_t \ge \xi > 0$; their Lemma~6 gives
$|q_t - p_t| \le \chi\delta$; and their Condition~1 is the mean contraction
$\int \|\varphi(S,W) - \varphi(S',W)\|\,\dd Q(W) \le \alpha \|S - S'\|$ with $\alpha < 1$.

\paragraph{FRS's Proposition 2.} Under Assumptions 1--11 and Condition 1, there are constants
$B$ and $L$ such that, for all $T$,
\[
  \bigl| L_j(y^T) - L(y^T) \bigr| < \Bigl(T B \chi + \frac{L}{1-\alpha}\Bigr)\delta .
\]

\section{Where the proof fails}

FRS split the gap as $|L_j - L| \le |\int F\,\dd\mu_j^* - \int F\,\dd\mu^*| +
|\int (F_j - F)\,\dd\mu_j^*|$, with $F = \int\prod_t p_t\,\dd Q$ and $F_j = \int \prod_t q_t
\,\dd Q$. Their bound (14) on the second term is $T B\chi\delta$. But the telescoping identity
\[
  \prod_{t=1}^T q_t - \prod_{t=1}^T p_t
  = \sum_{t=1}^T \Bigl(\prod_{s<t} q_s\Bigr)(q_t - p_t)\Bigl(\prod_{s>t} p_s\Bigr)
\]
places $T-1$ density factors beside each error $q_t - p_t$. With $M$ a bound on the densities the
honest bound is $T M^{T-1}\chi\delta$, and $M^{T-1}$ is not a constant: it grows exponentially
when $M > 1$ and vanishes exponentially when $M < 1$. This is exactly the factor
$(\sigma\sqrt{2\pi})^{-(T-1)}$ in Ackerberg, Geweke and Hahn's sharp bound for the Gaussian
location model. The step cannot be repaired with any $T$-free constant:

\begin{proposition}\label{prop:gap}
Let $M > 1$ and $\varepsilon > 0$. There is no $B$ such that
$|(M+\varepsilon)^T - M^T| \le T B\varepsilon$ for every $T \in \mathbb{N}$.
\end{proposition}
\formal{not\_exists\_linear\_bound\_prod\_sub}

\begin{proof}
By Bernoulli's inequality $(M+\varepsilon)^{n+1} - M^{n+1} \ge (n+1)\varepsilon M^n$, so the bound
would force $M^n \le B$ for all $n$.
\end{proof}

Two further steps let constants depend on $T$ without comment. The Lipschitz constant of
$F$ in $S_0$, used in (12), is ``different for each $t$, but since $t$ is finite, we can set a
global $L$''; the proposition is then stated ``for all $T$''. And Lemma~6's $\chi$ is asserted
for all $t$ ``since the model is stationary'', although the recovered $(v_t, s_t, w_{2t})$ depend
on a $t$-fold composition of the transition. We return to both in Section~\ref{sec:hyp}.

\section{The corrected proposition}

The accumulation of error that FRS describe is real, but it is \emph{multiplicative}: it holds
for the log likelihood, or equivalently the likelihood ratio, and not for the level.

\begin{assumption}[Uniform per-period error]\label{h:err}
$\xi \le p_t(s,\omega)$, $\xi \le q_t(s,\omega)$ and $|q_t(s,\omega) - p_t(s,\omega)| \le
\varepsilon$ for all $t, s, \omega$, with $\xi > 0$. In FRS's notation $\varepsilon = \chi\delta$.
\end{assumption}

\begin{assumption}[Log-Lipschitz in the initial state]\label{h:lip}
There are $\ell_t \ge 0$ with $p_t(s',\omega) \le e^{\ell_t d(s,s')} p_t(s,\omega)$ for all
$s, s', \omega$. Write $K = \sum_{t=1}^T \ell_t$. Only the \emph{exact} densities need this.
\end{assumption}

\begin{assumption}[Coupled initial states]\label{h:coup}
There is a probability measure $P$ on a space $X$ and maps $a, b : X \to S$ whose laws are
$\mu^*$ and $\mu_j^*$, with $d(a(x), b(x)) \le \rho$ for $P$-almost every $x$.
\end{assumption}

\begin{proposition}[Proposition 2$'$]\label{prop:main}
Under Hypotheses~\ref{h:err}--\ref{h:coup}, and integrability of the likelihood integrands,
\[
  e^{-D} L(y^T) \le L_j(y^T) \le e^{D} L(y^T), \qquad D = \frac{T\varepsilon}{\xi} + K\rho,
\]
equivalently $\bigl|\log L_j(y^T) - \log L(y^T)\bigr| \le D$.
\end{proposition}
\formal{likelihood\_le\_exp\_mul, abs\_log\_likelihood\_sub\_le}

\begin{proof}
\emph{Step 1 (one factor).} If $\xi \le p$ and $|q - p| \le \varepsilon$ then
$q \le p + \varepsilon \le p(1 + \varepsilon/\xi) \le e^{\varepsilon/\xi} p$.
\formal{le\_exp\_div\_mul}

\emph{Step 2 (the product; this replaces FRS's (14)).} Multiplying Step 1 over $t$,
$\prod_t q_t \le e^{T\varepsilon/\xi} \prod_t p_t$ pointwise in $(s,\omega)$.
\formal{prod\_le\_exp\_card\_mul\_prod}

\emph{Step 3 (moving the initial state).} Multiplying Hypothesis~\ref{h:lip} over $t$,
$\prod_t p_t(s',\omega) \le e^{K d(s,s')} \prod_t p_t(s,\omega)$.

\emph{Step 4 (integrating; this replaces FRS's (12)).} Combining Steps 2 and 3 and integrating
over $\omega$, $\int \prod_t q_t(s',\cdot)\,\dd Q \le e^{T\varepsilon/\xi + K d(s,s')}
\int \prod_t p_t(s,\cdot)\,\dd Q$ for all $s, s'$. Put $s = a(x)$, $s' = b(x)$, use
$d(a(x),b(x)) \le \rho$ almost surely, and integrate against $P$: $L_j \le e^D L$.
\formal{lik\_le\_exp\_mul\_lik}

For the reverse inequality, take the two moves in the other order: first move the initial state
from $a(x)$ to $b(x)$ \emph{within the exact model} (Step 3), then switch to the computed model at
$b(x)$ (Step 2, with the roles of $p$ and $q$ exchanged, which uses $\xi \le q_t$). This is why
Hypothesis~\ref{h:lip} is needed only for the exact densities.
\formal{lik\_le\_exp\_mul\_lik\_rev}
\noindent Since the densities are at least $\xi$ and $Q, P$ are probability measures,
$L, L_j \ge \xi^T > 0$, and taking logs gives the second form.
\end{proof}

\begin{remark}[Sharpness in $T$]
With every exact density equal to $M$ and every computed one to $M + \varepsilon$,
$\log L_j - \log L = T \log(1 + \varepsilon/M)$ exactly. The bound cannot grow more slowly than
linearly in $T$. In Ackerberg, Geweke and Hahn's location model restricted to a compact support,
taking every $y_t$ at the same point $\gamma + c$ gives
$\log L_j - \log L = T(c\delta - \delta^2/2)/\sigma^2$, again linear in $T$.
\end{remark}
\formal{log\_prod\_div\_prod\_const}

\section{The hypotheses from primitives}\label{sec:hyp}

Hypothesis~\ref{h:err} is FRS's Assumptions 4 and 8 and the conclusion of their Lemma~6, stated
uniformly in $t$. Hypotheses~\ref{h:lip} and~\ref{h:coup} replace their Condition~1 and the
Santos--Peralta-Alva bound used in (12). We derive all three from primitives on the model, with
proofs in the appendices. Two of FRS's assumptions are strengthened.

\begin{itemize}[leftmargin=1.5em]
\item \emph{Pathwise contraction.} Condition~1 is strengthened from contraction in mean to
contraction for every shock: $d(\varphi(s,w), \varphi(s',w)) \le \alpha\, d(s,s')$ for all $w$.
\item \emph{No $W_2$.} We take $\dim V_t = \dim Y_t$, so that the measurement shock alone is
recovered from the data. The state path $s_{t+1} = \varphi(s_t, \omega_t)$ is then a function of
the initial state and the shock path only, and the conditional density of the $t$-th observation
is $f_t(s_{t})$ for a function $f_t$ of the current state: the density of $V$ at
$g^{-1}(y_t; s)$ times the Jacobian.
\end{itemize}

\begin{theorem}[Proposition 2$'$ from primitives]\label{thm:prim}
Let $S$ be a compact metric space and $W$ a measurable space carrying the shock law $\nu$. Let
$\varphi, \varphi_j : S \times W \to S$ be jointly measurable and continuous in the state, with
$\varphi$ a pathwise $\alpha$-contraction ($0 \le \alpha < 1$) and
$d(\varphi_j(s,w), \varphi(s,w)) \le \delta$. Suppose each chain has a unique invariant
distribution, $\mu^*$ and $\mu_j^*$. Let the measurable per-state densities $f_t, f_{j,t}$ take
values in $[\xi, M]$ with $\xi > 0$, satisfy $|f_{j,t} - f_t| \le \eta$, and let each $f_t$ be
$c$-Lipschitz. Then for every sample length $T$ and every law $Q$ of the shock path,
\[
  \bigl|\log L_j(y^T) - \log L(y^T)\bigr| \le
  \frac{T}{\xi}\Bigl(\eta + \frac{c\,\delta}{1-\alpha}\Bigr) + \frac{c\,\alpha}{\xi}
  \frac{\delta}{(1-\alpha)^2}.
\]
\end{theorem}
\formal{abs\_log\_likelihood\_sub\_le\_of\_primitives}

\noindent Measurability and integrability of all the integrands are derived, not assumed. The
three ingredients:
\begin{itemize}[leftmargin=1.5em]
\item \emph{The coupling} (Appendix~\ref{app:coupling}). Drive both chains with the same shocks.
The band $\{(s,s') : d(s',s) \le \delta/(1-\alpha)\}$ is compact and invariant under the joint
chain, so by Krylov--Bogolyubov the joint chain has an invariant probability measure on it. Its
marginals are invariant for each chain, hence equal to $\mu^*$ and $\mu_j^*$.
\item \emph{Lemma~6 uniformly in $t$} (Appendix~\ref{app:prim}). The computed and exact state
paths from a common initial state stay within $\delta/(1-\alpha)$ of each other, so
$|q_t - p_t| \le \eta + c\delta/(1-\alpha)$ for every $t$.
\item \emph{Lipschitz in the initial state} (Appendix~\ref{app:prim}). The exact chain forgets its
initial state at rate $\alpha^t$, so $\ell_t = (c\alpha/\xi)\alpha^t$ and
$K \le c\alpha/(\xi(1-\alpha))$ for every $T$.
\end{itemize}

When $W_2$ is not empty, the recovered shocks depend on the data and on the whole past path, and
Hypotheses~\ref{h:err} and~\ref{h:lip} become a stability condition on the inverted filter. It is
not implied by FRS's Condition~1, and we have not found a primitive condition that delivers it.
The nearest results are on forgetting of the initial condition in hidden Markov models
\citep{BickelRitovRyden1998,DoucMoulinesRyden2004} and on uniform-in-time approximation of
nonlinear filters \citep{LeGlandOudjane2004}. They assume a transition density bounded above and
below, which does not fit FRS's setting, where the state moves by a deterministic map of the
shocks.

The general form with geometric decay, without the no-$W_2$ structure, is:

\begin{corollary}[Constants uniform in $T$]\label{cor:uniform}
Under Hypotheses~\ref{h:err}--\ref{h:coup}, if $\ell_t \le c\alpha^t$ and
$\rho = \delta/(1-\alpha)$, then for every $T$
\[
  \bigl|\log L_j(y^T) - \log L(y^T)\bigr| \le \frac{T\varepsilon}{\xi}
  + \frac{c\,\delta}{(1-\alpha)^2},
  \qquad
  \frac{1}{T}\bigl|\log L_j(y^T) - \log L(y^T)\bigr| \le
  \frac{\varepsilon}{\xi} + \frac{c\,\delta}{T(1-\alpha)^2}.
\]
\end{corollary}
\formal{abs\_log\_likelihood\_sub\_le\_of\_geometric}

This is FRS's bound, $T\cdot\mathrm{const}\cdot\delta + \mathrm{const}\cdot\delta/(1-\alpha)$,
with the level replaced by the log.

\section{Posteriors}\label{sec:post}

Suppose Proposition~2$'$ holds at every parameter value $\gamma$ with a common $D$, so that
$e^{-D}L(\gamma) \le L_j(\gamma) \le e^{D}L(\gamma)$ for all $\gamma$. Write
$\pi(\gamma \mid y^T) = L(\gamma)/Z$ and $\pi_j(\gamma\mid y^T) = L_j(\gamma)/Z_j$ for the
posterior densities with respect to a prior $\pi_0$, with $Z = \int L\,\dd\pi_0$ and
$Z_j = \int L_j\,\dd\pi_0$.

\begin{proposition}[Posteriors]\label{prop:post}
If $L, L_j > 0$ are $\pi_0$-integrable and $Z > 0$, then
\begin{enumerate}[label=(\roman*),leftmargin=2em]
\item $Z_j \le e^{D} Z$ and $Z \le e^{D} Z_j$;
\item $L_j(\gamma')/L_j(\gamma) \le e^{2D}\, L(\gamma')/L(\gamma)$ for all $\gamma, \gamma'$, so
Metropolis--Hastings acceptance ratios are off by at most a factor $e^{2D}$;
\item $\pi_j(\gamma\mid y^T) \le e^{2D}\,\pi(\gamma\mid y^T)$, and symmetrically;
\item $\int |\pi_j - \pi|\,\dd\pi_0 \le e^{2D} - 1$, so the total variation distance between the
posteriors is at most $(e^{2D}-1)/2$;
\item for measurable $h$ with $|h| \le B$, the posterior expectations satisfy
$|\mathbb{E}_j[h \mid y^T] - \mathbb{E}[h\mid y^T]| \le B\,(e^{2D}-1)$.
\end{enumerate}
\end{proposition}
\formal{evidence\_le\_exp\_mul, ratio\_le\_exp\_mul\_ratio, postDens\_le\_exp\_mul,
integral\_abs\_postDens\_sub\_le, abs\_integral\_postDens\_sub\_le}

Part (v) replaces FRS's Corollary~4, which asserts convergence $E_j(h\mid y^T) \to E(h\mid y^T)$
without a rate, with an explicit bound. The closest result in the literature is the
approximation theorem of \citet[Theorem~4.9 of arXiv:1302.6989v4]{DashtiStuart2017}, building on
\citet{CotterDashtiStuart2010}: if the negative log-likelihoods satisfy
$|\Phi - \Phi_N| \le M_2(\|u\|)\,\psi(N)$, then the Hellinger distance between the posteriors is
at most $C\psi(N)$ for $N$ large. Proposition~2$'$ supplies exactly that hypothesis, with
$\Phi = -\log L$ and $\psi = \delta$ (and $M_2$ proportional to $T$). Proposition~\ref{prop:post}
is the version with explicit constants that holds for every $\delta$; the proofs are in
Appendix~\ref{app:post}.

\section{What survives of FRS's conclusions}

\begin{itemize}[leftmargin=1.5em]
\item \emph{Error accumulates with $T$.} True of the log likelihood, at exactly the linear rate
FRS describe. In relative terms $|L_j/L - 1| \le e^{D} - 1 \approx D$ while $D$ is small, so
FRS's message ``the error is of order $T\delta$'' is correct for the \emph{relative} error of the
likelihood. Their statement for the absolute error is false.
\item \emph{Likelihood ratios, Metropolis--Hastings, posteriors (FRS Corollaries 1--4).} All hold
with explicit constants $e^{\pm 2D}$ (Proposition~\ref{prop:post}).
\item \emph{Estimation.} The average log likelihood $\frac1T\log L_j$ is within $O(\delta)$ of
$\frac1T\log L$, uniformly in $T$ (Theorem~\ref{thm:prim}, Corollary~\ref{cor:uniform}). This is
the object \citet{AGH2009} identify as relevant for the maximum likelihood estimator, and it
agrees with their conclusion that approximation error of order $\delta$ moves the estimator's
probability limit by $O(\delta)$, not by more as $T$ grows. A sup bound on the objective alone
gives only $O(\sqrt{\delta})$ for the maximiser under quadratic identification; their $O(\delta)$
rate uses closeness of derivatives as well.
\end{itemize}

What does not survive is FRS's conclusion that linearisation ``is due to fail as the sample size
grows'' because the likelihood levels diverge. The levels of both likelihoods go to zero or
infinity like $M^T$ anyway. What grows is the log-likelihood gap, and it grows no faster than the
log likelihood itself, so the average stays $O(\delta)$. And the posterior bounds of
Proposition~\ref{prop:post} degrade as $e^{2D}$ with $D \propto T\delta$: for Bayesian inference
FRS's warning stands, in the corrected form that $T\delta$ should be small.

\appendix

\section{The coupling of the invariant distributions}\label{app:coupling}

\begin{lemma}[Krylov--Bogolyubov for a random map]\label{lem:kb}
Let $X$ be a nonempty compact metric space, $\nu$ a probability measure on a measurable space $W$,
and $F : X \times W \to X$ jointly measurable with $F(\cdot, w)$ continuous for each $w$. Then the
chain $X_{t+1} = F(X_t, W_{t+1})$, $W_t \sim \nu$ i.i.d., has an invariant probability measure:
there is $\mu$ with $(\mu \otimes \nu) \circ F^{-1} = \mu$.
\end{lemma}
\formal{RandomMap.exists\_isInvariant}

\begin{proof}
Write $T\mu = (\mu\otimes\nu)\circ F^{-1}$ and $(Mh)(x) = \int h(F(x,w))\,\nu(\dd w)$ for bounded
continuous $h$. By change of variables and Fubini, $\int h\,\dd(T\mu) = \int Mh\,\dd\mu$
(\lean{integral\_push}). By dominated convergence ($h$ is bounded and $F(\cdot,w)$ is continuous),
$Mh$ is continuous (\lean{continuous\_markovFn}): this is the Feller property, and with the
duality it makes $T$ weakly continuous. Start from a Dirac measure $\mu_0$ and form the Cesàro
averages $\bar\mu_n = \frac{1}{n+1}\sum_{k=0}^n T^k\mu_0$. The difference
$\int Mh\,\dd\bar\mu_n - \int h\,\dd\bar\mu_n$ telescopes to
$\frac{1}{n+1}(\int h\,\dd T^{n+1}\mu_0 - \int h\,\dd\mu_0)$, which is at most $2\|h\|/(n+1)$
(\lean{integral\_cesaro\_sub}). By Prokhorov's theorem the probability measures on compact $X$ form
a compact space, so $(\bar\mu_n)$ has a cluster point $\mu$. By weak continuity,
$\int Mh\,\dd\mu = \int h\,\dd\mu$ for every bounded continuous $h$, i.e.
$\int h\,\dd(T\mu) = \int h\,\dd\mu$, and so $T\mu = \mu$.
\end{proof}

\noindent The repository already proved this for the income fluctuation problem
(\lean{exists\_isStationary} in \lean{Distribution/Stationary}), where the shock is finite; the
general version
differs only in the duality and the Feller step.

\begin{theorem}[The coupling]\label{thm:coupling}
Let $S$ be a nonempty compact metric space, $\varphi, \varphi_j$ jointly measurable and continuous
in the state, $\varphi$ a pathwise $\alpha$-contraction, $d(\varphi_j(s,w), \varphi(s,w)) \le
\delta$, $\delta \ge 0$. If $\mu^*$ is the unique invariant distribution of $\varphi$ and $\mu_j^*$
the unique invariant distribution of $\varphi_j$, there is a probability measure $P$ on the band
$B_\rho = \{(s,s') : d(s',s) \le \rho\}$, $\rho = \delta/(1-\alpha)$, whose first and second
marginals are $\mu^*$ and $\mu_j^*$.
\end{theorem}
\formal{exists\_coupling}

\begin{proof}
$B_\rho$ is closed in the compact space $S\times S$, so compact, and nonempty since it contains
the diagonal. Define the joint map $G((s,s'), w) = (\varphi(s,w), \varphi_j(s',w))$. It maps
$B_\rho$ into itself: $d(\varphi_j(s',w), \varphi(s,w)) \le d(\varphi_j(s',w),\varphi(s',w)) +
d(\varphi(s',w),\varphi(s,w)) \le \delta + \alpha\rho = \rho$ (\lean{dist\_step\_le\_of\_le}). It
is jointly measurable and continuous in the state pair (\lean{jointChain}). By
Lemma~\ref{lem:kb} there is an invariant probability measure $P$ on $B_\rho$. Taking the first
marginal commutes with one step of the chain, because the first coordinate of $G$ is $\varphi$ of
the first coordinate: $(P_1\otimes\nu)\circ\varphi^{-1} = ((P\otimes\nu)\circ G^{-1})_1 = P_1$
(\lean{push\_map\_marginal}). So $P_1$ is invariant for $\varphi$, and $P_1 = \mu^*$ by
uniqueness. The same argument gives $P_2 = \mu_j^*$.
\end{proof}

\noindent The theorem needs uniqueness of the computed chain's invariant distribution, which FRS
do not assume. Without it the conclusion holds for \emph{some} invariant distribution of the
computed chain, the one the coupling selects.

\section{The densities from primitives}\label{app:prim}

Let $s_t(s_0,\omega)$ be the exact state path ($s_0$ given, $s_{t+1} = \varphi(s_t, \omega_t)$)
and $s_{j,t}$ the computed one (\lean{statePath}). The likelihood's factors are
$p_t(s_0,\omega) = f_t(s_{t+1}(s_0,\omega))$ and $q_t = f_{j,t}(s_{j,t+1})$ (\lean{densPath}).

\begin{lemma}[Paths]\label{lem:paths}
For all $s_0, s_0', \omega, t$:
(a) $d(s_{j,t}(s_0,\omega), s_t(s_0,\omega)) \le \delta/(1-\alpha)$;
(b) $d(s_t(s_0,\omega), s_t(s_0',\omega)) \le \alpha^t d(s_0, s_0')$.
\end{lemma}
\formal{dist\_statePath\_le, dist\_statePath\_exact\_le}

\begin{proof}
(a) With $d_t = d(s_{j,t}, s_t)$, one step gives $d_{t+1} \le \alpha d_t + \delta$
(\lean{dist\_step\_le}), so by induction $d_t \le \alpha^t d_0 + \delta/(1-\alpha)$
(\lean{dist\_iterate\_le}), and $d_0 = 0$. (b) Induction on $t$ using the contraction.
\end{proof}

\begin{lemma}[Lemma 6, uniformly in $t$]\label{lem:unif}
$|q_t(s_0,\omega) - p_t(s_0,\omega)| \le \eta + c\,\delta/(1-\alpha)$ for all $t, s_0, \omega$.
\end{lemma}
\formal{abs\_densPath\_sub\_le}

\begin{proof}
With $x = s_{t+1}$ and $y = s_{j,t+1}$,
$|f_{j,t}(y) - f_t(x)| \le |f_{j,t}(y) - f_t(y)| + |f_t(y) - f_t(x)| \le \eta + c\,d(y,x)$, and
$d(y,x) \le \delta/(1-\alpha)$ by Lemma~\ref{lem:paths}(a).
\end{proof}

\begin{lemma}[Log-Lipschitz in the initial state]\label{lem:loglip}
$p_t(s_0', \omega) \le \exp\bigl((c\alpha/\xi)\,\alpha^t\, d(s_0, s_0')\bigr)\, p_t(s_0, \omega)$.
\end{lemma}
\formal{densPath\_le\_exp\_mul}

\begin{proof}
With $x = s_{t+1}(s_0,\omega)$ and $x' = s_{t+1}(s_0',\omega)$,
$|f_t(x') - f_t(x)| \le c\,d(x,x') \le c\,\alpha^{t+1} d(s_0,s_0')$ by Lemma~\ref{lem:paths}(b).
Since $f_t(x) \ge \xi$, the one-factor step of Proposition~\ref{prop:main} turns this additive
bound into the multiplicative bound $f_t(x') \le e^{c\alpha^{t+1}d(s_0,s_0')/\xi} f_t(x)$.
\end{proof}

\begin{proof}[Proof of Theorem~\ref{thm:prim}]
Theorem~\ref{thm:coupling} gives the coupling $P$ on the band, with $\rho = \delta/(1-\alpha)$.
Lemma~\ref{lem:unif} is Hypothesis~\ref{h:err} with $\varepsilon = \eta + c\delta/(1-\alpha)$.
Lemma~\ref{lem:loglip} is Hypothesis~\ref{h:lip} with $\ell_t = (c\alpha/\xi)\alpha^t$, and
$\sum_{t<T}\alpha^t \le 1/(1-\alpha)$ (\lean{sum\_le\_of\_geometric}). Corollary~\ref{cor:uniform}
then bounds the log-likelihood gap for the likelihoods written against $P$. Those equal the
likelihoods against $\mu^*$ and $\mu_j^*$ by change of variables along the marginal maps.

Measurability is derived. The state path is jointly measurable in $(s_0, \omega)$ by induction
(\lean{measurable\_statePath}), so the conditional likelihood $\prod_t f_t(s_{t+1})$ is jointly
measurable. By Fubini its integral over $\omega$ is measurable in $s_0$
(\lean{integrable\_densPath}). The bounds $0 \le \prod_t f_t \le M^T$ give integrability against
every probability measure.
\end{proof}

\section{Proofs for the posteriors}\label{app:post}

\begin{proof}[Proof of Proposition~\ref{prop:post}]
(i) Integrate $L_j \le e^D L$ against the prior. (ii) $L_j(\gamma') \le e^D L(\gamma')$ and
$1/L_j(\gamma) \le e^D/L(\gamma)$. (iii) Divide $L_j(\gamma) \le e^D L(\gamma)$ by
$Z_j \ge e^{-D} Z$. (iv) Write $a = \pi_j(\gamma\mid y^T)$, $b = \pi(\gamma\mid y^T)$, both
positive. By (iii) applied both ways, $a \le e^{2D} b$ and $b \le e^{2D} a$. Hence
$a - b \le (e^{2D}-1)\,b$, and $b - a \le b - e^{-2D} b = (1 - e^{-2D})\,b \le (e^{2D}-1)\,b$,
so $|a - b| \le (e^{2D}-1)\,b$ (\lean{abs\_postDens\_sub\_le}). Integrate, using
$\int b\,\dd\pi_0 = 1$. (v) $|\int h(a - b)\,\dd\pi_0| \le B \int|a-b|\,\dd\pi_0$, then (iv).
\end{proof}

\section{The formal development}

\begin{center}
\small
\begin{tabular}{@{}ll@{}}
\toprule
Result & Lean name \\
\midrule
\multicolumn{2}{@{}l}{\emph{\lean{Econometrics/ApproxLikelihood}}} \\
One factor, additive to multiplicative & \lean{le\_exp\_div\_mul} \\
Corrected product step (14) & \lean{prod\_le\_exp\_card\_mul\_prod} \\
FRS's step (14) cannot hold & \lean{not\_exists\_linear\_bound\_prod\_sub} \\
Sharpness in $T$ & \lean{log\_prod\_div\_prod\_const} \\
Coupled chain: one step, band, path & \lean{dist\_step\_le\_of\_le}, \lean{dist\_iterate\_le} \\
Proposition 2$'$, both directions & \lean{lik\_le\_exp\_mul\_lik}, \lean{\ldots\_rev} \\
Proposition 2$'$, log form & \lean{abs\_log\_likelihood\_sub\_le} \\
Corollary~\ref{cor:uniform} & \lean{abs\_log\_likelihood\_sub\_le\_of\_geometric} \\
\midrule
\multicolumn{2}{@{}l}{\emph{\lean{Distribution/RandomMap}}} \\
Lemma~\ref{lem:kb} (Krylov--Bogolyubov) & \lean{RandomMap.exists\_isInvariant} \\
\midrule
\multicolumn{2}{@{}l}{\emph{\lean{Econometrics/ApproxLikelihoodModel}}} \\
Theorem~\ref{thm:coupling} (coupling) & \lean{exists\_coupling} \\
Lemma~\ref{lem:paths} & \lean{dist\_statePath\_le}, \lean{dist\_statePath\_exact\_le} \\
Lemma~\ref{lem:unif} & \lean{abs\_densPath\_sub\_le} \\
Lemma~\ref{lem:loglip} & \lean{densPath\_le\_exp\_mul} \\
Theorem~\ref{thm:prim} & \lean{abs\_log\_likelihood\_sub\_le\_of\_primitives} \\
\midrule
\multicolumn{2}{@{}l}{\emph{\lean{Econometrics/ApproxPosterior}}} \\
Proposition~\ref{prop:post} (i)--(iii) & \lean{evidence\_le\_exp\_mul}, \lean{postDens\_le\_exp\_mul} \\
Proposition~\ref{prop:post} (iv) & \lean{integral\_abs\_postDens\_sub\_le} \\
Proposition~\ref{prop:post} (v) & \lean{abs\_integral\_postDens\_sub\_le} \\
\bottomrule
\end{tabular}
\end{center}

\noindent Paths are relative to \lean{LeanEconomics/}. The library builds without \lean{sorry}, and every
result above depends only on Lean's three standard axioms (\lean{propext},
\lean{Classical.choice}, \lean{Quot.sound}).

What remains informal: the reduction of FRS's measurement equation to per-state densities $f_t$
(the density of $V$ at $g^{-1}(y_t;s)$ times the Jacobian) and the bound $\eta = O(\delta)$ on
their error at a fixed state, which is FRS's Lemma~6 argument without the time dimension; and the
case with $W_2$ present, discussed in Section~\ref{sec:hyp}.

\bibliographystyle{plainnat}
\bibliography{refs}

\end{document}
